\section{Future Scope}

\subsection{Completed Milestones}

Since the initial submission of this progress report, several components that
were previously identified as future work have been successfully implemented
and deployed.

\subsubsection{Mobile-Friendly Web Application}
A fully functional web-based interface has been developed and deployed,
enabling users to upload environmental images and receive real-time AQI
predictions in simplified categorical classes --- \textit{Healthy},
\textit{Moderate}, and \textit{Unhealthy} --- without requiring any local
setup. The application is cross-device accessible and optimized for both
desktop and mobile browsers.

\subsubsection{Cloud Deployment and Backend Infrastructure}
The backend has been deployed on cloud infrastructure, ensuring scalability
and real-time processing. The complete backend source code is publicly
available at:
\begin{center}
  \url{https://github.com/ASR222/AQI_project}
\end{center}
The repository includes the trained model weights, FastAPI-based inference
endpoints, and deployment configuration for reproducibility.

\subsection{Current Future Scope}

\subsubsection{Video-Based AQI Classification}
The next major technical milestone is extending the current image-based
classification pipeline to support continuous video streams. Rather than
classifying individual static frames, the proposed system will perform
temporal aggregation across consecutive frames to produce smoother and more
reliable AQI estimates in real time. This will involve integrating lightweight
temporal models or sliding-window inference over live camera feeds, making the
system suitable for deployment on fixed surveillance infrastructure, traffic
cameras, and drone footage.

\subsubsection{Government and Institutional Tie-Ups}
We are actively pursuing collaborations with relevant government bodies and
environmental agencies to integrate our image-based AQI estimation system as
a supplementary tool within existing air quality monitoring networks. Such
partnerships would enable access to geographically distributed camera
infrastructure, ground-truth sensor data for model validation, and pathways
toward official recognition of image-based monitoring as a credible
supplementary methodology. The goal is to bridge the gap between academic
research and actionable environmental policy through institutional adoption.
